%!TEX TS-program = xe\hbox{\LaTeX}
\documentclass[10pt]{article}

\usepackage[a4paper,margin=1.55cm]{geometry}
\usepackage{fontspec}
\usepackage{paracol}
\usepackage{titlesec}
\usepackage{hyperref}
\usepackage{xcolor}
\usepackage{enumitem}

\definecolor{royalblue}{HTML}{003399}

\setmainfont[
    Path = ./,
    UprightFont = *-Regular,
    ItalicFont  = *-Italic,
    BoldFont    = *-Smbd,
    BoldItalicFont = *-SmbdItalic
]{ArnoPro}

\titleformat{\section}
  {\large\bfseries\color{royalblue}}
  {}
  {0pt}
  {}
  [\vspace{-3pt}\color{royalblue}\titlerule]

\setlength{\parindent}{0pt}
\setlength{\parskip}{2.2pt}
\setlength{\emergencystretch}{1em}
\setlist[itemize]{left=7pt, itemsep=-0.5pt, topsep=2pt, parsep=0pt}

\newcommand{\role}[2]{%
  \item {\fontsize{11.4pt}{13pt}\selectfont\textbf{#1}}\hfill\textit{#2}\\[-1pt]
}

\hypersetup{
    colorlinks=true,
    urlcolor=royalblue,
    linkcolor=royalblue
}

\begin{document}

\noindent
\begin{minipage}[t]{0.55\textwidth}
    {\huge \textbf{Srikanth Mohankumar}}\\[4pt]
    {\large\raggedright\textit{Engineering Lead | SaaS \& Automation Platforms}}
\end{minipage}
\begin{minipage}[t]{0.43\textwidth}
    \raggedleft
    Tamil Nadu, India\\
    +91\,8526629711\\
    \href{mailto:srikanthmohankumar@gmail.com}{srikanthmohankumar@gmail.com}\\
    \href{https://www.linkedin.com/in/srikanth-mohankumar}{linkedin.com/in/srikanth-mohankumar}\\
    \href{https://github.com/Srikanth-Mohankumar}{github.com/Srikanth-Mohankumar}\\
    \href{https://ctan.org/author/mohankumar}{ctan.org/author/mohankumar}
\end{minipage}

\vspace{-8pt}

\columnratio{0.25,0.70}
\setlength{\columnsep}{0.65cm}
\begin{paracol}{2}

\section*{Skills}
\textbf{SaaS \& Architecture}\\
Cloud-native design, APIs, scalable backends, product thinking\\[-8pt]

\textbf{AI \& LLMs}\\
AI agents, RAG, vector stores, NLP, prompt engineering\\[-8pt]

\textbf{Cloud / DevOps}\\
AWS, GCP, Docker, Linux, CI and CD\\[-8pt]

\textbf{Backend \& Data}\\
Python, Node.js, FastAPI, Redis, MongoDB, REST APIs\\[-8pt]

\textbf{Publishing Systems}\\
LaTeX, LuaTeX, TeX4ht, XML, XSLT, PDF, accessibility\\[-8pt]

\textbf{Leadership}\\
Mentoring, Scrum, RCA, delivery coordination

\section*{Certifications}
\textbf{GenAI/LLM Programme}\\
TNQ Tech, 40 hours: LLMs, agents, embeddings, automation.

\textbf{AI Credentials}\\
LangGraph agents, Anthropic AI Fluency, Claude Code — Mar 2026.

\section*{Conferences}
\textbf{TUG 2025}\\
In-person, Kerala, India — structure, language, accessibility, and TeX.\\
\href{https://tug.org/tug2025}{tug.org/tug2025}

\textbf{TUG 2026}\\
Remote participation — Calgary programme on typography, publishing, tagged PDF.\\
\href{https://tug.org/tug2026}{tug.org/tug2026}

\section*{Education}
\textbf{B.E. Electrical \& Electronics}\\
Kingston Engineering College, Vellore, Tamil Nadu.

\switchcolumn

\section*{Summary}
Engineering Lead and SaaS Product Builder with 8+ years in \textbf{publishing technology}, building cloud-native automated typesetting platforms, AI-assisted content transformation systems, and workflow tools. Promoted to \textbf{Engineering Lead} at TNQ Tech in \textbf{April 2026}. Strong product mindset across feasibility analysis, user-story refinement, sprint flow, RCA, mentoring, and stakeholder-facing prototypes.

\section*{Career Evolution \& Experience}
\begin{itemize}[leftmargin=*, label={\textbullet}]

\role{Engineering Lead — Product, Prototyping \& SaaS Automation}{Apr 2026 -- Present}
\textit{TNQ Tech, Chennai, Tamil Nadu, India}
\begin{itemize}
  \item Lead platform engineering for automated typesetting, SaaS delivery, GenAI demos, and stakeholder-facing prototypes
  \item Coordinate product, production, QA, and engineering teams on ambiguous enterprise publishing workflow problems
  \item Mentor engineers and production technologists; improve predictability through Scrum, Kanban, RCA, and workflow design
\end{itemize}

\role{Software Engineer — Product \& Prototyping}{Dec 2022 -- Mar 2026}
\textit{TNQ Tech, Chennai, Tamil Nadu, India}
\begin{itemize}
  \item Core contributor to \textbf{NeoPage}, a cloud-based automated typesetting SaaS platform reducing licensed-tool dependency
  \item Achieved \textbf{2x performance improvement} via on-the-fly PDF pagination (10--30s vs.\ 60+s)
  \item Built distributed content transformation services using Python, Node.js, FastAPI, Docker, AWS Lambda, and MongoDB
  \item Built LLM demos and GenAI showcases to modernize legacy publishing pipelines
\end{itemize}

\role{Team Lead}{Apr 2022 -- Nov 2022}
\textit{Lapiz Digital Services, Ranipet, Tamil Nadu, India}
\begin{itemize}
  \item Led a 35+ member production team across SLA-heavy publishing projects; managed delivery, QC, RCA, mentoring, workload planning, and automation-oriented improvements
\end{itemize}

\role{Template Developer}{Apr 2021 -- Mar 2022}
\textit{Integra Software Services, Puducherry, India}
\begin{itemize}
  \item Developed LaTeX journal/book templates, pagination logic, reusable style components, and publisher-specific production workflows
\end{itemize}

\role{LaTeX Production \& Automation}{Sep 2017 -- Mar 2021}
\textit{Lapiz Digital Services, Ranipet, Tamil Nadu, India}
\begin{itemize}
  \item Progressed from Trainee to Technical Associate; built LaTeX macros, scripts, reusable tooling, and workflow standards
  \item Achieved \textbf{4x productivity improvement}; handled complex pagination, revision cycles, and junior mentoring
\end{itemize}

\end{itemize}

\section*{Role Impact}
\textbf{Platform delivery:} SaaS publishing workflows, automated typesetting, and stakeholder-facing prototypes.\\
\textbf{Operational lift:} workflow design, RCA, sprint rhythm, mentoring, and reusable production tooling.\\
\textbf{Document engineering:} TeX, XML, PDF generation, browser pagination, and accessibility-aware publishing systems.

\end{paracol}

\newpage
\section*{Selected Projects}
\textbf{NeoView — PDF Viewer for Technical Workflows} \hfill \textit{2026}\\
Python, PySide6, PyMuPDF · \href{https://srikanth-mohankumar.github.io/neoview/}{Website}\\
Technical PDF viewer for font inspection, measurement, crop/export, annotations, and review-heavy publishing workflows.

\textbf{authpack — LuaLaTeX Author Management} \hfill \textit{2025}\\
LuaLaTeX · \href{https://github.com/Srikanth-Mohankumar/latex-authpack}{GitHub}\\
Author-affiliation tooling with ORCID support and flexible scholarly publishing layouts.

\textbf{enumsub — CTAN LaTeX Package} \hfill \textit{2025}\\
\href{https://ctan.org/pkg/enumsub}{CTAN} · \href{https://ctan.org/author/mohankumar}{Contributor Profile}\\
CTAN package for inline sublists and consistent enumeration alignment.

\textbf{GenAI-Powered Editorial Assistant} \hfill \textit{2025}\\
Streamlit, GenAI, NLP · \href{https://github.com/Srikanth-Mohankumar/email-sentiment-analysis}{GitHub}\\
Capstone project that reads, interprets, and drafts responses to emails using NLP and sentiment analysis with a Streamlit UI and containerized backend.

\textbf{Browser-Powered Pagination} \hfill \textit{2024}\\
TeX, Lua, Emscripten, WebAssembly\\
Custom browser pagination engine that runs TeX/Lua through WebAssembly for full pagination functionality without a desktop publishing dependency.

\textbf{Support ChatBot} \hfill \textit{2023}\\
Level 1 support automation for production-team query handling. Helped production users answer repeated workflow queries faster and keep support knowledge closer to daily operations.

\section*{Professional Focus}
I work best on systems where publishing rules, human review, and automation must meet: SaaS editorial platforms, PDF workflows, LaTeX infrastructure, AI-assisted content operations, and practical tools for production teams. Current interests include tagged and accessible PDF, browser-native pagination, agentic editorial assistance, and scalable document transformation pipelines.

\section*{Project Impact Highlights}
\textbf{Publishing automation:} reusable systems for pagination, content transformation, template logic, and review support. \textbf{Performance:} improved PDF responsiveness from 60+ seconds to 10--30 seconds. \textbf{AI prototyping:} GenAI demos around sentiment, response drafting, retrieval, and agentic support. \textbf{Team enablement:} support tooling, reusable LaTeX components, and technical handoff material.

\section*{Engineering Toolkit}
\textbf{Application:} Python, Node.js, FastAPI, Streamlit, REST APIs, Redis, MongoDB.\\
\textbf{Cloud \& Delivery:} AWS Lambda, Docker, Linux, CI/CD workflows, feasibility analysis, MVP scoping.\\
\textbf{Publishing:} LaTeX, LuaLaTeX, LuaTeX, TeX4ht, XML, XSLT, PDF generation, accessibility tagging concepts.\\
\textbf{AI:} LLM agents, LangGraph learning, RAG, embeddings, vector search, NLP, prompt engineering.

\section*{Community \& Learning}
CTAN contributor through \textbf{enumsub}; continued open-source publishing-tool experiments through GitHub. Attended \textbf{TUG 2025} in person in Kerala; following \textbf{TUG 2026} remotely for accessible PDF, typography, and TeX ecosystem updates. Completed TNQ Tech GenAI/LLM programme and March 2026 AI credentials.

\section*{Working Style}
Start with real production constraints: exception cases, review points, file formats, and quality gates. Turn unclear ideas into demos, feasibility notes, user stories, and MVP slices teams can evaluate together. Current direction: accessible/tagged PDF automation, browser-native TeX workflows, and agentic tools that support editors and production teams.

\end{document}
